% Luc Destombe --- licence CC BY-NC-SA 4.0
% https://creativecommons.org/licenses/by-nc-sa/4.0/deed.fr
% Fichier autonome : compiler avec pdflatex, deux fois (signets, nombre de pages).
\documentclass[11pt,a4paper]{article}
\makeatletter
\newif\ifeleve
\elevefalse
\elevetrue

\newif\iflycee@palatino
\lycee@palatinotrue

\RequirePackage[utf8]{inputenc}
\RequirePackage[T1]{fontenc}
\RequirePackage[french]{babel}
\RequirePackage{etoolbox}

\newcommand{\ShorthandsOff}{\@ifundefined{shorthandoff}{}{\shorthandoff{;:!?}}}
\newcommand{\ShorthandsOn}{\@ifundefined{shorthandon}{}{\shorthandon{;:!?}}}
\ShorthandsOff
\AtBeginDocument{\ShorthandsOn}
\AtBeginEnvironment{tikzpicture}{\ShorthandsOff}

\RequirePackage{geometry}
\RequirePackage{fancyhdr}
\RequirePackage{lastpage}
\RequirePackage{multicol}
\RequirePackage{array}
\RequirePackage{enumitem}
\RequirePackage{xcolor}
\RequirePackage{needspace}

\RequirePackage{amsmath,amssymb}
\RequirePackage{mathpazo}
\DeclareMathSizes{10.95}{10.95}{8}{5.5}
\begingroup\catcode`\_=\active \catcode`\^=\active
\gdef_#1{\sb{\scriptscriptstyle #1}}
\gdef^#1{\sp{\scriptscriptstyle\lycee@moinsepais #1}}
\endgroup
\newcommand\lycee@moins{\mathbin{%
  \setbox\z@\hbox{$\m@th\scriptscriptstyle\lycee@moinsorig$}%
  \dimen@=.55\wd\z@ \advance\dimen@-.5pt
  \hbox to.75\wd\z@{\hss$\m@th\scriptscriptstyle\vcenter{\hbox{%
    \tikz\draw[line width=.5pt,line cap=round](0,0)--(\the\dimen@,0);}}$\hss}}}
\begingroup\lccode`\~=`\- \lowercase{\endgroup\let~}\lycee@moins
\newcommand\lycee@moinsepais{\mathcode`\-="8000 }
\AtBeginDocument{\mathchardef\lycee@moinsorig=\mathcode`\-\relax}
\AtBeginDocument{\catcode`\_=12 \mathcode`\_="8000
  \catcode`\^=12 \mathcode`\^="8000 }
\newcommand\lycee@exposants{%
  \fontdimen13\textfont2=.48\fontdimen6\textfont2
  \fontdimen14\textfont2=.48\fontdimen6\textfont2
  \fontdimen15\textfont2=.48\fontdimen6\textfont2
  \fontdimen13\scriptfont2=.48\fontdimen6\scriptfont2
  \fontdimen14\scriptfont2=.48\fontdimen6\scriptfont2
  \fontdimen15\scriptfont2=.48\fontdimen6\scriptfont2}
\every@math@size\expandafter{\the\every@math@size\lycee@exposants}
\newlength{\retraitcalculs}
\setlength{\retraitcalculs}{2em}

\RequirePackage{tikz}
\usetikzlibrary{arrows.meta,calc,babel,shapes.misc}
\RequirePackage[most]{tcolorbox}
\tcbuselibrary{skins,breakable}

\definecolor{coulDef}{HTML}{1F4E79}
\definecolor{coulProp}{HTML}{7030A0}
\definecolor{coulRem}{HTML}{C55A11}
\definecolor{coulEx}{HTML}{2E7D32}
\definecolor{coulExo}{HTML}{B03A2E}
\definecolor{coulPart}{HTML}{1F4E79}
\definecolor{coulMeth}{HTML}{1B6E4F}
\definecolor{coulCourbe}{HTML}{0B5ED7}
\definecolor{coulCourbeB}{HTML}{C00000}
\definecolor{coulCourbeC}{HTML}{2E7D32}
\definecolor{coulGrille}{HTML}{8C8C8C}
\definecolor{coulRep}{HTML}{1B6E4F}
\definecolor{coulBar}{HTML}{2E7D32}

\newcommand{\R}{\mathbb{R}}
\newcommand{\Rs}{\mathbb{R}^{*}}
\renewcommand{\le}{\leqslant}
\renewcommand{\ge}{\geqslant}
\newcommand{\vect}[1]{\overrightarrow{#1}}
\newcommand{\norme}[1]{\left\lVert #1 \right\rVert}
\newcommand{\intervalle}[2]{\left[#1\,;#2\right]}
\RequirePackage{cancel}
\renewcommand{\CancelColor}{\color{coulExo}}
\newcommand{\fois}[1]{\times{\color{coulExo}#1}}
\newcommand{\barre}[1]{\cancel{#1}}

\tikzset{grille/.style={coulGrille,line width=0.4pt}}
\newcommand{\quadrillage}[4]{\draw[grille,step=1] (#1,#2) grid (#3,#4);}
\tikzset{
  croix/.style={cross out,draw,line width=0.6pt,inner sep=0pt,minimum size=4pt},
  croixdroite/.style={croix,rotate=45,line width=0.8pt},
}
\newcommand{\pt}[3]{\path (#1) node[croix]{} node[#2,font=\small,inner sep=2.5pt]{$#3$};}

\newcommand{\lycee@repere}[4]{%
  \draw[grille,step=1] (#1,#2) grid (#3,#4);
  \draw[-{Stealth[length=2mm]},thick] (#1,0)--({#3+0.4},0) node[below left=-1pt]{$x$};
  \draw[-{Stealth[length=2mm]},thick] (0,#2)--(0,{#4+0.4}) node[below left=-1pt]{$y$};
  \node[below left,font=\scriptsize,inner sep=1.5pt] at (0,0) {$0$};}
\newcommand{\lycee@gradx}[1]{\foreach \k/\l in {#1}{%
  \draw (\k,0.07)--(\k,-0.07) node[below,font=\scriptsize,inner sep=1.5pt]{$\l$};}}
\newcommand{\lycee@grady}[1]{\foreach \k/\l in {#1}{%
  \draw (0.07,\k)--(-0.07,\k) node[left,font=\scriptsize,inner sep=1.5pt]{$\l$};}}
\newcommand{\lycee@intff}[2]{\left[#1\,;#2\right]}
\newcommand{\lycee@intfo}[2]{\left[#1\,;#2\right[}
\newcommand{\lycee@intof}[2]{\left]#1\,;#2\right]}
\newcommand{\lycee@intoo}[2]{\left]#1\,;#2\right[}
\newcommand{\lycee@ens}[1]{\left\{#1\right\}}
\AtBeginDocument{%
  \providecommand{\repere}{\lycee@repere}%
  \providecommand{\gradx}{\lycee@gradx}%
  \providecommand{\grady}{\lycee@grady}%
  \providecommand{\Cf}{\mathcal{C}\sb{\scriptscriptstyle f}}%
  \providecommand{\Cg}{\mathcal{C}\sb{\scriptscriptstyle g}}%
  \providecommand{\intff}{\lycee@intff}%
  \providecommand{\intfo}{\lycee@intfo}%
  \providecommand{\intof}{\lycee@intof}%
  \providecommand{\intoo}{\lycee@intoo}%
  \providecommand{\ens}{\lycee@ens}}

\newlist{questions}{enumerate}{3}
\setlist[questions,1]{label=\arabic*\textdegree),leftmargin=*,itemsep=2pt,topsep=3pt}
\setlist[questions,2]{label=\alph*),leftmargin=*,itemsep=1pt,topsep=2pt}
\setlist[questions,3]{label=\roman*),leftmargin=*,itemsep=1pt,topsep=1pt}
\setlist[itemize]{label=\textbullet,leftmargin=*,itemsep=2pt,topsep=3pt}

\newcommand{\besoinplace}[1]{\par\penalty\@M\begingroup
  \dimen@=#1\relax\dimen@ii=\pagegoal\advance\dimen@ii-\pagetotal
  \ifdim\dimen@>\dimen@ii\ifdim\dimen@ii>\z@\vfil\fi\break\fi\endgroup}

\newcommand{\lycee@pied}{\thepage/\pageref{LastPage}}
\newcommand{\entete}[2]{%
  \pagestyle{fancy}\fancyhf{}%
  \fancyhead[L]{\footnotesize\color{coulPart}\textbf{#1}}%
  \fancyhead[R]{\footnotesize\color{coulPart}\textbf{#2}}%
  \fancyfoot[C]{\footnotesize\lycee@pied}%
  \renewcommand{\headrulewidth}{0.4pt}%
  \renewcommand{\headrule}{\hbox to\headwidth{%
    \color{coulPart!40}\leaders\hrule height \headrulewidth\hfill}}}

\RequirePackage{ccicons}
\newcommand{\lycee@licence}{{\scriptsize\color{gray}%
  \copyright\ Luc Destombe --- \ccbyncsa\ CC BY-NC-SA 4.0}}
\newif\iflycee@licence \lycee@licencetrue
\newcommand{\sanslicence}{\lycee@licencefalse}
\AtBeginDocument{\iflycee@licence\AddToHookNext{shipout/foreground}{%
  \put(0,0){\rlap{\kern\dimexpr1in+\hoffset+\oddsidemargin+\textwidth\relax
    \llap{\raisebox{-\dimexpr1in+\voffset+\topmargin+\headheight+\headsep
      +\textheight+\footskip\relax}{\lycee@licence}}}}}\fi}

\newcommand{\formulesserrees}{%
  \setlength{\abovedisplayskip}{3pt}\setlength{\belowdisplayskip}{3pt}%
  \setlength{\abovedisplayshortskip}{2pt}\setlength{\belowdisplayshortskip}{3pt}}

\setlength{\parindent}{0pt}

\RequirePackage[implicit=false,hidelinks,pdfencoding=auto,psdextra,bookmarksopen,
  bookmarksopenlevel=1,pdfcreator={},pdfproducer={}]{hyperref}
\RequirePackage{bookmark}
\hypersetup{pdfauthor={Luc Destombe},pdfkeywords={CC BY-NC-SA 4.0}}
\newcounter{lycee@signet}
\newcommand{\signet}[3][]{\stepcounter{lycee@signet}%
  \edef\lycee@dest{\if\relax\detokenize{#1}\relax
    signet.\arabic{lycee@signet}\else#1\fi}%
  \hypertarget{\lycee@dest}{}\bookmark[level=#2,dest=\lycee@dest]{#3}}
\pdfstringdefDisableCommands{%
  \def\R{R}\def\vect#1{#1}\def\Cf{Cf}\def\Cg{Cg}\def\fois#1{×#1}%
  \def\barre#1{#1}\def\intervalle#1#2{[#1;#2]}\def\intff#1#2{[#1;#2]}%
  \def\textsuperscript#1{#1}\def\dfrac#1#2{#1/#2}\def\frac#1#2{#1/#2}%
  \def\vec#1{#1⃗}}%
\geometry{top=2cm,bottom=1cm,left=1cm,right=1cm,headheight=14pt,footskip=0.6cm}

\RequirePackage{pgfplots}
\pgfplotsset{compat=1.18}
\RequirePackage{tkz-tab}

\newcounter{partiecours}
\renewcommand{\thepartiecours}{\Roman{partiecours}}
\newcommand{\partiecours}[1]{%
  \refstepcounter{partiecours}%
  \Needspace*{8\baselineskip}%
  \bigskip
  \noindent\begin{tcolorbox}[enhanced,colback=coulPart,colframe=coulPart,
      boxrule=0pt,arc=2pt,left=8pt,right=8pt,top=4pt,bottom=4pt,
      before skip=6pt,after skip=10pt]
    \color{white}\bfseries\large \thepartiecours{} --- #1\signet{1}{\thepartiecours{} --- #1}
  \end{tcolorbox}%
  \addcontentsline{toc}{section}{\thepartiecours{} --- #1}%
}

\newtcolorbox{boitecours}[3][]{%
  enhanced,breakable,
  colback=white,colframe=#2,coltitle=white,
  fonttitle=\bfseries,
  attach boxed title to top left={xshift=8pt,yshift=-\tcboxedtitleheight/2},
  boxed title style={colback=#2,arc=2pt,boxrule=0pt},
  boxrule=0.9pt,arc=2pt,
  left=8pt,right=8pt,top=8pt,bottom=6pt,
  title={#3},#1}

\newcounter{methode}
\newenvironment{definitionbox}[1][Définition]%
  {\begin{boitecours}[phantom={\signet{2}{#1}}]{coulDef}{#1}}{\end{boitecours}}
\newenvironment{proprietebox}[1][Propriété]%
  {\begin{boitecours}[phantom={\signet{2}{#1}}]{coulProp}{#1}}{\end{boitecours}}
\newenvironment{methodebox}[1][Méthode]{\stepcounter{methode}%
  \begin{boitecours}[phantom={\signet[methode-\arabic{methode}]{2}{#1}}]{coulMeth}{#1}}%
  {\end{boitecours}}

\newcommand{\etape}[1]{\par\smallskip\textbf{\color{coulMeth}#1}\par\nobreak\smallskip}
\newcommand{\enonce}[1]{\emph{#1}\par\smallskip}
\newcommand{\reponse}[1]{\par\smallskip
  {\footnotesize\itshape\color{black!55}Réponses : #1}}
\newcommand{\demo}[1]{\textbf{\color{coulMeth}Démonstration.} #1}
\newcommand{\afaire}[1]{%
  \par\smallskip
  \noindent{\small\color{coulExo}$\blacktriangleright$\ \textbf{À faire :}
    fiche d'exercices du chapitre, #1.}\par\smallskip}

\newenvironment{calculs}{\setlength{\jot}{5pt}\@fleqntrue\@mathmargin\retraitcalculs
  \setlength{\abovedisplayskip}{4pt}\setlength{\belowdisplayskip}{4pt}%
  \setlength{\abovedisplayshortskip}{4pt}\setlength{\belowdisplayshortskip}{4pt}%
  \csname alignat*\endcsname{2}}%
  {\csname endalignat*\endcsname}
\newcommand{\rem}[1]{\text{\small\itshape\color{black!60}#1}}
\newcommand{\explique}[2]{\par\smallskip\noindent
  \begin{minipage}[t]{0.57\linewidth}\raggedright #1\end{minipage}\hfill
  \begin{minipage}[t]{0.40\linewidth}\raggedright\leavevmode
    {\small\itshape\color{black!60}#2}\end{minipage}\par}
\newcommand{\cas}[1]{\par\smallskip\noindent\textbf{\color{coulExo}#1}\nobreak}
\newcommand{\calc}[1]{\par\smallskip\textbf{\color{coulExo}#1}\par\nobreak}

\newtcolorbox{remarquebox}[1][Remarque]{%
  enhanced,breakable,colback=white,colframe=coulRem,
  boxrule=0pt,leftrule=2.5pt,arc=0pt,outer arc=0pt,
  left=8pt,right=6pt,top=5pt,bottom=5pt,
  before upper={\textbf{\color{coulRem}#1}\par\smallskip}}

\newtcolorbox{exemplebox}[1][Exemple]{%
  enhanced,breakable,colback=white,colframe=coulEx,
  boxrule=0pt,leftrule=2.5pt,arc=0pt,outer arc=0pt,
  left=8pt,right=6pt,top=5pt,bottom=5pt,
  before upper={\textbf{\color{coulEx}#1}\par\smallskip}}

\pgfplotsset{
  repere/.style={
    axis lines=middle,
    axis line style={-{Stealth[length=2.4mm]},thick},
    every axis x label/.style={at={(ticklabel* cs:1.0)},anchor=west},
    every axis y label/.style={at={(ticklabel* cs:1.0)},anchor=south},
    label style={font=\footnotesize},
    tick label style={font=\scriptsize},
    grid=both,
    major grid style={grille},
    minor tick num=0,
    tick align=inside,
    clip mode=individual,
  },
}
\tikzset{
  courbe/.style={coulCourbe,line width=1.1pt,smooth},
  courbeB/.style={coulCourbeB,line width=1.1pt,smooth},
  courbeC/.style={coulCourbeC,line width=1.1pt,smooth},
}

\newlist{etapes}{enumerate}{1}
\setlist[etapes]{label=\textbf{\arabic*.},leftmargin=*,itemsep=1pt,topsep=2pt}

\newlist{expressions}{enumerate}{1}
\setlist[expressions]{label={\Alph*\ $=$},leftmargin=*,itemsep=3pt,topsep=3pt}

\newcounter{exo}
\renewcommand{\theexo}{\ifnum\value{exo}<10 0\fi\arabic{exo}}
\newtcolorbox{exercice}[1][]{%
  enhanced,breakable,
  colback=white,colframe=coulExo,
  boxrule=0.9pt,arc=2pt,
  left=8pt,right=8pt,top=8pt,bottom=6pt,
  coltitle=white,fonttitle=\bfseries,
  attach boxed title to top left={xshift=8pt,yshift=-\tcboxedtitleheight/2},
  boxed title style={colback=coulExo,arc=2pt,boxrule=0pt},
  phantom={\signet{2}{Exercice \arabic{exo}}},
  title={Exercice \theexo\ifx\relax#1\relax\else{} --- #1\fi}}
\newenvironment{exo}[1][\relax]{\refstepcounter{exo}\begin{exercice}[#1]}{\end{exercice}}

  \RequirePackage{comment}
  \excludecomment{correction}
  \renewcommand{\reponse}[1]{}
  \newcommand{\autableau}{\hfill{\footnotesize\itshape\color{black!45}(au tableau)}}
  \newcommand{\etapeexemples}{%
    \par\smallskip\textbf{\color{coulMeth}Exemples}\autableau\par\nobreak\smallskip}
  \newcommand{\etapetableau}[1]{%
    \par\smallskip\textbf{\color{coulMeth}#1}\autableau\par\nobreak\smallskip}
  \newtcolorbox{exempletableau}[1][Exemple]{%
    enhanced,breakable,colback=white,colframe=coulEx,
    boxrule=0pt,leftrule=2.5pt,arc=0pt,outer arc=0pt,
    left=8pt,right=6pt,top=4pt,bottom=4pt,
    before skip=5pt,after skip=5pt,
    before upper={\textbf{\color{coulEx}#1}\autableau\par\smallskip}}

\newcommand{\coursEntete}{}
\newcommand{\coursNiveau}{}
\newcommand{\coursTitre}{}
\newcommand{\coursSurTitre}{}
\newcommand{\coursSousTitre}{}

\newcommand{\enteteCours}[1]{\renewcommand{\coursEntete}{#1}}
\newcommand{\chapitrecours}[2]{\enteteCours{Chapitre #1 --- #2}}
\newcommand{\niveaucours}[1]{\renewcommand{\coursNiveau}{#1}}
\newcommand{\titrecours}[3][]{%
  \renewcommand{\coursTitre}{#2}%
  \renewcommand{\coursSurTitre}{#1}%
  \renewcommand{\coursSousTitre}{#3}}

\pagestyle{fancy}
\fancyhf{}
\fancyhead[L]{\footnotesize\color{coulPart}\textbf{\coursEntete}}
\fancyhead[R]{\footnotesize\color{coulPart}\textbf{\coursNiveau}}
\fancyfoot[C]{\footnotesize\thepage}
\renewcommand{\headrulewidth}{0.4pt}
\renewcommand{\headrule}{\hbox to\headwidth{\color{coulPart!40}\leaders\hrule height \headrulewidth\hfill}}

\setlength{\parskip}{2pt}

\AtBeginDocument{%
  \begin{center}
    \begin{tcolorbox}[enhanced,width=0.72\linewidth,colback=white,colframe=coulPart,
        boxrule=1.6pt,arc=3pt,halign=center,top=10pt,bottom=10pt]
      \ifx\coursSurTitre\empty
        {\Huge\bfseries\color{coulPart} \coursTitre}\\[4pt]
      \else
        {\Huge\bfseries\color{coulPart} \coursTitre}\\[3pt]
        {\Large\bfseries\color{coulPart} \coursSurTitre}\\[5pt]
      \fi
      {\normalsize \coursSousTitre}%
      \\[2pt]{\small\itshape Version élève}
    \end{tcolorbox}
  \end{center}%
}
\makeatother

\chapitrecours{3}{Le second degré}
\niveaucours{1\textsuperscript{re} spécialité mathématiques}
\titrecours{LE SECOND DEGRÉ}{Cours --- Première spécialité}

\begin{document}

\providecommand{\panneauattention}{\tikz[baseline=-0.25em]{%
  \draw[red!75!black,line width=0.8pt,line join=round,fill=yellow!25]
    (-0.5em,-0.35em)--(0.5em,-0.35em)--(0,0.55em)--cycle;
  \node[font=\bfseries\scriptsize,inner sep=0pt] at (0,-0.02em) {!};}}

\begin{remarquebox}[Comment utiliser ce cours]
Chaque partie commence par ce qu'il faut \textbf{savoir} (définitions, théorèmes), puis
vient ce qu'il faut \textbf{savoir faire} : une méthode, avec des
\textbf{exemples} faits en classe, puis des exercices
\og\textbf{À vous de jouer}\fg{} à chercher seul.

\smallskip
Sur cette feuille, les exemples n'ont que leur \textbf{énoncé} : la résolution, marquée
\emph{(au tableau)}, s'écrit dans le cahier pendant le cours. Les réponses des « À vous
de jouer » ne sont pas données.

\end{remarquebox}

\partiecours{La forme canonique du trinôme}

\begin{definitionbox}[Définition 1 --- Trinôme du second degré]
Un \textbf{trinôme du second degré} (ou \textbf{trinôme}) est une expression de la forme
\[
  P(x)=ax^{2}+bx+c \qquad\text{où $a$, $b$, $c$ sont des réels, avec } a\ne 0.
\]
\end{definitionbox}

\begin{exemplebox}[Exemple]
$P(x)=3x^{2}-5x+2$ est un trinôme : $a=3$, $b=-5$ et $c=2$.
\end{exemplebox}

Un même trinôme peut s'écrire de plusieurs façons : forme développée, forme canonique et
forme factorisée (si elle existe). La \textbf{forme canonique} est celle qui montre s'il
se factorise et quel est son extremum. Pour l'obtenir, on fait apparaître une
\textbf{identité remarquable}.

\begin{exemplebox}[Exemple --- compléter un carré]
Soit $P(x)=x^{2}+8x+5$. On sait que $(x+4)^{2}=x^{2}+8x+16$ : les deux premiers termes
de $P(x)$ sont le début de ce carré, auquel il manque $16$.
\begin{calculs}
P(x) &= x^{2}+8x+16-16+5 &\qquad& \rem{on ajoute $16$ et on le retire aussitôt}\\
P(x) &= (x+4)^{2}-11 &\qquad& \rem{forme $(x-\alpha)^{2}+\beta$ avec $\alpha=-4$ et $\beta=-11$}
\end{calculs}
\end{exemplebox}

\begin{proprietebox}[Théorème 1 --- Forme canonique]
Tout trinôme $P(x)=ax^{2}+bx+c$ s'écrit, d'une seule façon,
\[
  \boxed{P(x)=a(x-\alpha)^{2}+\beta
  \qquad\text{avec}\qquad \alpha=-\dfrac{b}{2a} \quad\text{et}\quad \beta=P(\alpha)}
\]
\end{proprietebox}

Dans l'exemple, $a=1$ et $b=8$ :
\begin{calculs}
\alpha &= -\dfrac{8}{2}=-4 &&\\
\beta  &= P(-4)=16-32+5 &&\\
\beta  &= -11 &\qquad& \rem{on retrouve la forme $(x+4)^{2}-11$}
\end{calculs}

\besoinplace{16\baselineskip}
\begin{remarquebox}[Démonstration (non exigible)]
C'est le même calcul que dans l'exemple, avec des lettres.
\begin{calculs}
P(x) &= a\left(x^{2}+\dfrac{b}{a}x\right)+c
  &\qquad& \rem{on met $a\ne 0$ en facteur dans les deux premiers termes}\\
P(x) &= a\left[\left(x+\dfrac{b}{2a}\right)^{2}-\dfrac{b^{2}}{4a^{2}}\right]+c
  &\qquad& \rem{$x^{2}+\frac{b}{a}x$ est le début de $\left(x+\frac{b}{2a}\right)^{2}$}\\
P(x) &= a\left(x+\dfrac{b}{2a}\right)^{2}-\dfrac{b^{2}}{4a}+c
  &\qquad& \rem{on distribue $a$}
\end{calculs}
On reconnaît $a(x-\alpha)^{2}+\beta$ avec $\alpha=-\dfrac{b}{2a}$. Pour $\beta$, inutile
de retenir une formule : en $x=\alpha$ le carré est nul, donc $P(\alpha)=\beta$.
\end{remarquebox}

\begin{methodebox}[Méthode 1 --- Mettre un trinôme sous forme canonique]
\etapeexemples
Écrire chaque trinôme sous forme canonique.
\begin{questions}[label=\alph*)]
  \item $P(x)=x^{2}-10x+16$
  \item $Q(x)=2x^{2}+10x+7$ \quad ($a\ne 1$)
  \item $R(x)=-x^{2}+6x-11$ \quad ($a<0$)
\end{questions}

\etape{À vous de jouer}
Écrire chaque trinôme sous forme canonique, puis contrôler avec $\alpha$ et $\beta$ :
\[
  A(x)=x^{2}+12x+20 \qquad B(x)=2x^{2}-2x-3 \qquad C(x)=-3x^{2}+12x-14
\]

\end{methodebox}

\begin{remarquebox}[À quoi sert la forme canonique ?]
Dans $a(x-\alpha)^{2}+\beta$, on regarde les signes de $a$ et de $\beta$ :
\begin{itemize}
  \item \textbf{signes contraires} : on factorise avec $u^{2}-v^{2}=(u-v)(u+v)$ ;
  \item \textbf{même signe} ($\beta\ne 0$) : le trinôme ne s'annule jamais, il ne se
        factorise pas.
\end{itemize}
\end{remarquebox}

\begin{exemplebox}[Exemple --- factoriser grâce à la forme canonique]
$P(x)=x^{2}-10x+16=(x-5)^{2}-9$, avec $a=1>0$ et $\beta=-9<0$ :
\begin{calculs}
P(x) &= (x-5)^{2}-3^{2} &\qquad& \rem{forme $u^{2}-v^{2}$}\\
P(x) &= (x-5-3)(x-5+3) &&\\
P(x) &= (x-8)(x-2) &&
\end{calculs}
\end{exemplebox}

\afaire{exercices 1 à 4}

\partiecours{Racines du trinôme et équations du second degré}

\begin{definitionbox}[Définition 2 --- Racines et discriminant]
Les \textbf{racines} d'un trinôme sont les solutions de l'équation $ax^{2}+bx+c=0$.

\smallskip
Le \textbf{discriminant} du trinôme est le nombre $\Delta=b^{2}-4ac$.
\end{definitionbox}

\begin{proprietebox}[Théorème 2 --- Nombre de racines]
\begin{itemize}
  \item Si $\Delta>0$ : deux racines
        $\;x_{1}=\dfrac{-b+\sqrt{\Delta}}{2a}\;$ et $\;x_{2}=\dfrac{-b-\sqrt{\Delta}}{2a}$.
  \item Si $\Delta=0$ : une seule racine (dite \textbf{double}) $\;x_{0}=-\dfrac{b}{2a}$.
  \item Si $\Delta<0$ : aucune racine réelle.
\end{itemize}
\end{proprietebox}

\begin{remarquebox}[Pourquoi ?]
En reprenant la démonstration de la partie I, la forme canonique s'écrit
\[
  P(x)=a\left[(x-\alpha)^{2}-\dfrac{\Delta}{4a^{2}}\right]
\]
Si $\Delta<0$, le crochet est la somme d'un carré et d'un nombre strictement positif : il
ne s'annule pas. Si $\Delta=0$, il reste $a(x-\alpha)^{2}$, nul seulement en $\alpha$. Si
$\Delta>0$, le crochet est une différence de deux carrés, qui se factorise et donne deux
racines.
\end{remarquebox}

\begin{methodebox}[Méthode 2 --- Résoudre une équation du second degré]
\etapeexemples
Résoudre dans $\R$ les équations suivantes.
\begin{questions}[label=\alph*)]
  \item $3x^{2}-4x-4=0$
  \item $5x^{2}-20x+20=0$
  \item $x^{2}-3x+5=0$
  \item $9x^{2}-16=0$ \quad et \quad $4x^{2}+12x+9=0$
\end{questions}

\etape{À vous de jouer}
Résoudre dans $\R$ :
\[
  x^{2}-9x+20=0 \qquad 3x^{2}+2x+1=0 \qquad 9x^{2}+6x+1=0 \qquad 5x^{2}+15x=0
\]

\end{methodebox}

\afaire{exercices 5 à 9}

\partiecours{Factorisation du trinôme}

Quand $\Delta>0$, la différence de deux carrés de la partie II se factorise ; ses deux
facteurs s'annulent en $x_{1}$ et en $x_{2}$. On obtient :

\begin{proprietebox}[Théorème 3 --- Forme factorisée]
\begin{itemize}
  \item Si $\Delta>0$ : $\;P(x)=a\,(x-x_{1})(x-x_{2})$.
  \item Si $\Delta=0$ : $\;P(x)=a\,(x-x_{0})^{2}$.
  \item Si $\Delta<0$ : le trinôme \textbf{ne se factorise pas}.
\end{itemize}
\end{proprietebox}

\begin{remarquebox}[\panneauattention\ Le facteur $a$]
On n'oublie pas $a$ devant les parenthèses. Contrôle rapide : le terme constant de
$a(x-x_{1})(x-x_{2})$ est $a\,x_{1}x_{2}$, qui doit valoir $c$.
\end{remarquebox}

\begin{methodebox}[Méthode 3 --- Factoriser un trinôme]
\etapeexemples
Factoriser, quand c'est possible, les trinômes suivants.
\begin{questions}[label=\alph*)]
  \item $P(x)=3x^{2}-4x-4$
  \item $Q(x)=5x^{2}-20x+20$
  \item $R(x)=x^{2}-3x+5$
  \item $T(x)=x^{2}+2x-15$
\end{questions}

\etape{À vous de jouer}
Factoriser quand c'est possible :
\[
  A(x)=x^{2}+x-12 \qquad B(x)=3x^{2}-7x+2 \qquad C(x)=x^{2}-2x+3 \qquad D(x)=-2x^{2}+3x+2
\]

\end{methodebox}

\afaire{exercices 10 à 12}

\partiecours{Somme et produit des racines}

On suppose que le trinôme a deux racines $x_{1}$ et $x_{2}$. On note $S=x_{1}+x_{2}$
leur \textbf{somme} et $P=x_{1}x_{2}$ leur \textbf{produit}. On développe la forme
factorisée :
\begin{calculs}
ax^{2}+bx+c &= a(x-x_{1})(x-x_{2}) &&\\
\color{black!60}(x-x_{1})(x-x_{2}) &\color{black!60}= x^{2}-(x_{1}+x_{2})x+x_{1}x_{2}
  &\qquad& \rem{calcul intermédiaire : on développe}\\
ax^{2}+bx+c &= a\left(x^{2}-Sx+P\right) &\qquad& \rem{avec $S=x_{1}+x_{2}$ et $P=x_{1}x_{2}$}\\
ax^{2}+bx+c &= ax^{2}-aSx+aP &&
\end{calculs}
En identifiant les coefficients : $b=-aS$ et $c=aP$.

\begin{proprietebox}[Théorème 4 --- Somme et produit des racines]
Si le trinôme $ax^{2}+bx+c$ a deux racines, leur somme $S$ et leur produit $P$ valent
\[
  \boxed{S=-\dfrac{b}{a} \qquad\text{et}\qquad P=\dfrac{c}{a}}
\]
En pratique, on met $a$ en facteur : $ax^{2}+bx+c=a\left(x^{2}-Sx+P\right)$.
\end{proprietebox}

\begin{exemplebox}[Exemple --- lire $S$ et $P$]
\[
  3x^{2}-4x-4=3\left(x^{2}-\dfrac{4}{3}x-\dfrac{4}{3}\right)
\]
donc $S=\dfrac{4}{3}$ et $P=-\dfrac{4}{3}$. Avec les racines $2$ et $-\dfrac{2}{3}$ de la
méthode 2 : $2-\dfrac{2}{3}=\dfrac{4}{3}$ et $2\times\left(-\dfrac{2}{3}\right)=-\dfrac{4}{3}$.
\end{exemplebox}

\begin{definitionbox}[Définition 3 --- Racine évidente]
Une \textbf{racine évidente} est une racine que l'on trouve de tête, en essayant des
valeurs simples : $1$, $-1$, $2$, $-2$\dots
\end{definitionbox}

\begin{methodebox}[Méthode 4 --- Trouver les racines sans discriminant]
\etapeexemples
Résoudre les équations suivantes sans calculer de discriminant.
\begin{questions}[label=\alph*)]
  \item $3x^{2}-8x+5=0$
  \item $4x^{2}+3x-1=0$
  \item $x^{2}-8x+12=0$
\end{questions}

\etape{À vous de jouer}
Trouver une racine évidente, puis l'autre racine avec la somme :
\[
  x^{2}+5x+4=0 \qquad 4x^{2}-7x+3=0 \qquad x^{2}+2x-8=0 \qquad -3x^{2}+2x+1=0
\]

\end{methodebox}

\afaire{exercices 13 à 18}

\partiecours{Signe du trinôme et inéquations du second degré}

\subsection*{Deux racines}
Si $\Delta>0$, avec $x_{1}<x_{2}$, le signe de $P(x)=a(x-x_{1})(x-x_{2})$ se lit dans un
tableau de signes :

\begin{center}
\begin{tikzpicture}
  \tkzTabInit[lgt=3.4,espcl=2.9]{$x$/1,$x-x_{1}$/0.9,$x-x_{2}$/0.9,$a(x-x_{1})(x-x_{2})$/1.1}%
    {$-\infty$,$x_{1}$,$x_{2}$,$+\infty$}
  \tkzTabLine{,-,z,+,t,+,}
  \tkzTabLine{,-,t,-,z,+,}
  \tkzTabLine{,\text{signe de } a,z,\text{signe de } {-a},z,\text{signe de } a,}
\end{tikzpicture}
\end{center}

Le trinôme a le signe de $a$ \textbf{à l'extérieur} des racines, le signe contraire
\textbf{entre} les racines.

\subsection*{Une seule racine ou aucune}
Si $\Delta=0$, $P(x)=a(x-x_{0})^{2}$ : un carré est positif, donc $P(x)$ a le signe de
$a$, et s'annule en $x_{0}$.

Si $\Delta<0$, $P(x)=a\left[(x-\alpha)^{2}-\dfrac{\Delta}{4a^{2}}\right]$ et le crochet est
strictement positif : $P(x)$ a toujours le signe de $a$.

\begin{proprietebox}[Théorème 5 --- Signe du trinôme]
Le signe de $P(x)$ dépend de $\Delta$ et du signe de $a$ :
\begin{center}
\begin{minipage}[t]{0.32\linewidth}\centering
{\small $\Delta>0$}\\[2pt]
\begin{tikzpicture}[scale=0.8]
  \tkzTabInit[lgt=1.2,espcl=1.2]{$x$/0.8,$P(x)$/0.8}{$-\infty$,$x_{1}$,$x_{2}$,$+\infty$}
  \tkzTabLine{,a,z,-a,z,a,}
\end{tikzpicture}
\end{minipage}\hfill
\begin{minipage}[t]{0.3\linewidth}\centering
{\small $\Delta=0$}\\[2pt]
\begin{tikzpicture}[scale=0.8]
  \tkzTabInit[lgt=1.2,espcl=1.4]{$x$/0.8,$P(x)$/0.8}{$-\infty$,$x_{0}$,$+\infty$}
  \tkzTabLine{,a,z,a,}
\end{tikzpicture}
\end{minipage}\hfill
\begin{minipage}[t]{0.3\linewidth}\centering
{\small $\Delta<0$}\\[2pt]
\begin{tikzpicture}[scale=0.8]
  \tkzTabInit[lgt=1.2,espcl=2.8]{$x$/0.8,$P(x)$/0.8}{$-\infty$,$+\infty$}
  \tkzTabLine{,a,}
\end{tikzpicture}
\end{minipage}
\end{center}
Dans chaque case, « $a$ » signifie « du signe de $a$ ».
\end{proprietebox}

\begin{exemplebox}[Exemple --- signe de $-4x^{2}+4x+3$]
\begin{calculs}
\Delta &= 4^{2}-4\times(-4)\times 3 &&\\
\Delta &= 64 &\qquad& \rem{deux racines}\\
x_{1} &= \dfrac{-4+8}{-8}=-\dfrac{1}{2} &&\\
x_{2} &= \dfrac{-4-8}{-8}=\dfrac{3}{2} &&
\end{calculs}
$a=-4<0$ : négatif à l'extérieur des racines, positif entre elles.
\begin{center}
\begin{tikzpicture}[scale=0.95]
  \tkzTabInit[lgt=2.8,espcl=2.4]{$x$/1,$-4x^{2}+4x+3$/1.2}%
    {$-\infty$,$-\frac{1}{2}$,$\frac{3}{2}$,$+\infty$}
  \tkzTabLine{,-,z,+,z,-,}
\end{tikzpicture}
\end{center}
\end{exemplebox}

\begin{methodebox}[Méthode 5 --- Résoudre une inéquation du second degré]
\etapeexemples
Résoudre dans $\R$ les inéquations suivantes.
\begin{questions}[label=\alph*)]
  \item $x^{2}-5x+4>0$
  \item $-x^{2}+2x+4>0$
  \item $x^{2}+2x+6>0$
  \item $3x^{2}-12x+12\le 0$
\end{questions}

\etape{À vous de jouer}
Résoudre :
\[
  x^{2}+2x-8<0 \qquad x^{2}-6x+5\ge 0 \qquad x^{2}+x+2<0 \qquad 2x^{2}+8x+8>0
\]

\end{methodebox}

\afaire{exercices 19 à 21}

\partiecours{Représentation graphique de la fonction trinôme}

Dans cette partie, $f(x)=ax^{2}+bx+c=a(x-\alpha)^{2}+\beta$ sur $\R$.

\begin{exemplebox}[Exemple --- l'extremum se lit sur la forme canonique]
Soit $f(x)=3(x+2)^{2}-1$. Pour tout réel $x$ :
\begin{calculs}
(x+2)^{2} &\ge 0 &\qquad& \rem{un carré est positif}\\
3(x+2)^{2}-1 &\ge -1 &\qquad& \rem{on multiplie par $3>0$, puis on retranche $1$}\\
f(x) &\ge -1 &&
\end{calculs}
Et $f(-2)=-1$ : $f$ a un \textbf{minimum}, égal à $-1$, atteint en $-2$.
\end{exemplebox}

\begin{definitionbox}[Définition 4 --- Parabole et sommet]
La courbe d'une fonction trinôme est une \textbf{parabole}. Son \textbf{sommet} est le
point $S(\alpha\,;\beta)$.
\end{definitionbox}

\begin{proprietebox}[Théorème 6 --- Orientation et extremum]
\begin{itemize}
  \item Si $a>0$ : parabole \textbf{tournée vers le haut} ; $f$ a un \textbf{minimum}
        $\beta$, atteint en $\alpha$.
  \item Si $a<0$ : parabole \textbf{tournée vers le bas} ; $f$ a un \textbf{maximum}
        $\beta$, atteint en $\alpha$.
\end{itemize}
La parabole est \textbf{symétrique} par rapport à la droite d'équation $x=\alpha$.
\end{proprietebox}

\begin{center}
\begin{minipage}[t]{0.47\linewidth}
\centering {\small $a>0$ : $f(x)=x^{2}+2x-3$}\\[2pt]
\begin{tikzpicture}
\begin{axis}[repere,
    x=0.55cm,y=0.55cm,
    xmin=-4.6,xmax=2.6,ymin=-5.2,ymax=4.8,
    xtick={-4,-3,-2,-1,1,2},ytick={-5,-4,-3,-2,-1,1,2,3,4},
    xlabel={$x$},ylabel={$y$}]
  \addplot[courbe,domain=-4.05:2.05,samples=100]{x^2+2*x-3};
  \addplot[only marks,mark=x,mark options={line width=0.6pt},mark size=2pt,black] coordinates {(-3,0)(1,0)(-1,-4)};
  \node[above right,font=\footnotesize] at (axis cs:-3,0) {$x_{2}$};
  \node[above left,font=\footnotesize] at (axis cs:1,0) {$x_{1}$};
  \node[below right,font=\scriptsize] at (axis cs:-1,-4) {$S(\alpha\,;\beta)$};
  \addplot[dashed,black!45] coordinates {(-1,-5.2)(-1,4.8)};
\end{axis}
\end{tikzpicture}
\end{minipage}\hfill
\begin{minipage}[t]{0.47\linewidth}
\centering {\small $a<0$ : $g(x)=-x^{2}-2x+3$}\\[2pt]
\begin{tikzpicture}
\begin{axis}[repere,
    x=0.55cm,y=0.55cm,
    xmin=-4.6,xmax=2.6,ymin=-4.8,ymax=5.2,
    xtick={-4,-3,-2,-1,1,2},ytick={-4,-3,-2,-1,1,2,3,4,5},
    xlabel={$x$},ylabel={$y$}]
  \addplot[courbe,domain=-4.05:2.05,samples=100]{-x^2-2*x+3};
  \addplot[only marks,mark=x,mark options={line width=0.6pt},mark size=2pt,black] coordinates {(-3,0)(1,0)(-1,4)};
  \node[above left,font=\footnotesize] at (axis cs:-3,0) {$x_{1}$};
  \node[above right,font=\footnotesize] at (axis cs:1,0) {$x_{2}$};
  \node[above right,font=\scriptsize] at (axis cs:-1,4) {$S(\alpha\,;\beta)$};
  \addplot[dashed,black!45] coordinates {(-1,-4.8)(-1,5.2)};
\end{axis}
\end{tikzpicture}
\end{minipage}
\end{center}

\begin{remarquebox}[Racines et parabole]
Les racines sont les abscisses des points où la parabole \textbf{coupe l'axe des
abscisses} : deux points si ${\Delta>0}$, un seul (le sommet) si ${\Delta=0}$, aucun si
${\Delta<0}$. Avec deux racines, $\alpha$ est leur \textbf{milieu} : sur les figures,
${\dfrac{-3+1}{2}=-1=\alpha}$.
\end{remarquebox}

\begin{methodebox}[Méthode 6 --- Tracer l'allure d'une parabole]
\etapeexemples
Tracer l'allure de la courbe de chaque fonction.
\begin{questions}[label=\alph*)]
  \item $f(x)=x^{2}+2x-3$
  \item $g(x)=-x^{2}-2x+3$
  \item $h(x)=x^{2}+4x+6$
\end{questions}

\etape{À vous de jouer}
Pour chaque fonction : orientation, sommet, racines éventuelles, puis allure de la courbe.
\[
  f_{1}(x)=x^{2}-4x+3 \qquad f_{2}(x)=-x^{2}+6x-5 \qquad f_{3}(x)=x^{2}-2x+5
\]

\end{methodebox}

\afaire{exercices 22 à 26}

\partiecours{Variations de la fonction trinôme}

\begin{proprietebox}[Théorème 7 --- Tableau de variations]
\begin{center}
\begin{minipage}{0.47\linewidth}\centering
{\small $a>0$ : minimum $\beta$ en $\alpha$}\\[4pt]
\begin{tikzpicture}[scale=0.95]
  \tkzTabInit[lgt=2,espcl=2.6]{$x$/1,$f(x)$/2}{$-\infty$,$\alpha$,$+\infty$}
  \tkzTabVar{+/{},-/$\beta$,+/{}}
\end{tikzpicture}
\end{minipage}\hfill
\begin{minipage}{0.47\linewidth}\centering
{\small $a<0$ : maximum $\beta$ en $\alpha$}\\[4pt]
\begin{tikzpicture}[scale=0.95]
  \tkzTabInit[lgt=2,espcl=2.6]{$x$/1,$f(x)$/2}{$-\infty$,$\alpha$,$+\infty$}
  \tkzTabVar{-/{},+/$\beta$,-/{}}
\end{tikzpicture}
\end{minipage}
\end{center}
avec $\alpha=-\dfrac{b}{2a}$ et $\beta=f(\alpha)$.
\end{proprietebox}

\begin{methodebox}[Méthode 7 --- Dresser le tableau de variations d'un trinôme]
\etapeexemples
Dresser le tableau de variations de chaque fonction, puis donner son extremum.
\begin{questions}[label=\alph*)]
  \item $f(x)=x^{2}+8x+10$
  \item $g(x)=-3x^{2}-6x-7$
  \item $h(x)=(x+1)(x-5)$
\end{questions}

\etape{À vous de jouer}
Dresser le tableau de variations et donner l'extremum :
\[
  f_{1}(x)=x^{2}+6x+4 \qquad f_{2}(x)=-2(x-3)^{2}+7 \qquad f_{3}(x)=3x^{2}-6x-1
\]

\end{methodebox}

\afaire{exercices 27 à 30}

\partiecours{Équations se ramenant au second degré}

Une équation qui contient des produits, des carrés ou des quotients devient souvent, après
calcul, une équation $ax^{2}+bx+c=0$.

\begin{methodebox}[Méthode 8 --- Résoudre une équation se ramenant au second degré]
\etapeexemples
Résoudre dans $\R$ les équations suivantes.
\begin{questions}[label=\alph*)]
  \item Un produit : $(x+3)(x-2)=14$.
  \item Deux carrés : $(3x-1)^{2}=(x+1)^{2}$.
  \item Un quotient : $\dfrac{6}{x}=x+1$.
  \item Une valeur à rejeter : $\dfrac{x^{2}-4x+4}{x-2}=3x-4$.
  \item Deux quotients : $\dfrac{1}{x-1}-\dfrac{2}{2x-1}=\dfrac{1}{3}$.
\end{questions}

\etape{À vous de jouer}
Résoudre dans $\R$ :
\[
  (x+2)(x-3)=6 \qquad (2x+1)^{2}=(x+3)^{2} \qquad \dfrac{8}{x}=x+2 \qquad
  \dfrac{2x^{2}+5x-3}{x+3}=x+1
\]

\end{methodebox}

\afaire{exercices 31 à 33}

\partiecours{Inéquations se ramenant au second degré}

\begin{methodebox}[Méthode 9 --- Résoudre une inéquation se ramenant au second degré]
\etapeexemples
Résoudre dans $\R$ les inéquations suivantes.
\begin{questions}[label=\alph*)]
  \item Un produit : $(x+4)(x-1)<3x+12$.
  \item Un quotient de trinômes : $\dfrac{x^{2}-x-2}{x^{2}+2x-3}\ge 0$.
  \item Deux quotients : $\dfrac{x+1}{x-2}<\dfrac{2}{x+1}$.
\end{questions}

\etape{À vous de jouer}
Résoudre :
\[
  (x+1)^{2}\ge(2x-1)^{2} \qquad \dfrac{x+3}{x-1}\ge 0 \qquad \dfrac{x^{2}-9}{x+1}\le 0
\]

\end{methodebox}

\begin{remarquebox}[Erreur à éviter]
Dans une inéquation, multiplier les deux membres par un dénominateur change le sens de
l'inégalité quand il est négatif. Comme on ne connaît pas son signe, on étudie le signe
d'un quotient.
\end{remarquebox}

\afaire{exercices 34 à 36}

\partiecours{Résoudre un problème du second degré}

\begin{methodebox}[Méthode 10 --- Résoudre un problème du second degré]
\etapeexemples
\begin{questions}[label=\alph*)]
  \item \textbf{Deux pages d'un livre.} On ouvre un livre au hasard. Le produit des
        numéros des deux pages visibles vaut $156$. Quels sont ces numéros ?
  \item \textbf{Un panneau solaire.} Un panneau solaire rectangulaire a un périmètre de
        $22$ dm et une aire de $28\text{ dm}^{2}$. Quelles sont ses dimensions ?
  \item \textbf{Une balle lancée d'une terrasse.} Sa hauteur au-dessus du sol, en
        mètres, $t$ secondes après le lancer, est $h(t)=-5t^{2}+10t+40$ pour $t\ge 0$.
        \textbf{1)} Quelle hauteur maximale atteint-elle, et quand ?
        \textbf{2)} Quand touche-t-elle le sol ?
  \item \textbf{Agrandir une terrasse.} Une terrasse rectangulaire mesure $6$ m sur $4$ m.
        On allonge ses deux côtés d'une même longueur $x$, et son aire passe à
        $48\text{ m}^{2}$. Trouver $x$.
\end{questions}

\etape{À vous de jouer}
\begin{questions}
  \item Deux élèves ont des numéros de casier consécutifs, dont le produit vaut $210$.
        Quels sont ces numéros ?
  \item Une table rectangulaire a un périmètre de $26$ dm et une aire de
        $42\text{ dm}^{2}$ : quelles sont ses dimensions ?
  \item Une fusée à eau est lancée depuis un balcon ; sa hauteur en mètres est
        $h(t)=-5t^{2}+30t+35$ pour $t\ge 0$. Quelle hauteur maximale atteint-elle ? Au
        bout de combien de temps retombe-t-elle au sol ?
\end{questions}

\end{methodebox}

\afaire{exercices 37 à 41}

\end{document}
